\Needspace{16\baselineskip}
\section*{AI use}\label{sec:ai-use}
The underlying research was carried out by autonomous rStation agents
using Gemini 4 Argon and GPT-6 Astra. The author launched the agents,
which conducted the research without human intervention. The sources for
the earlier bound are
\href{https://rstation.ai/problems/tenth_discriminant/papers/2/v12}{paper 2,
version 12} and
\href{https://rstation.ai/problems/tenth_discriminant/papers/22/v1}{paper 22,
version 1}. The source for the bound $331$ is
\href{https://rstation.ai/problems/tenth_discriminant/papers/24}{rStation paper 24}.
Its calculations extend
\href{https://rstation.ai/problems/tenth_discriminant/papers/2/v13}{paper 2,
version 13}, together with the rank argument developed in
\href{https://rstation.ai/problems/tenth_discriminant/papers/22/v4}{paper 22,
version 4}. The spectral estimates originate in
\href{https://rstation.ai/problems/tenth_discriminant/papers/16/v2}{paper 16,
version 2}. 

Separate agents using GPT-6 Astra and Claude Opus 5.5 checked the original
arguments independently of the discovering agents. GPT-6 Astra, Claude
Opus 5.5 and Gemini 4 Argon were then used in preparing the present
exposition, including mathematical consultation, revision of the outline
and prose, and review of the proof. GPT-6 Astra recomputed the six numerical comparisons for this paper
and used separate programs to check the rank inequalities, the
auxiliary-prime counts and the remaining discriminants. The source
calculation at $347$ included a four-factor composite correction;
the comparisons here use only the two- and three-factor corrections.
